\documentclass[11pt]{article}

\usepackage[spanish]{babel}
\usepackage[utf8]{inputenc}
\usepackage{amsmath, amssymb}
\usepackage{geometry,xcolor}
\geometry{margin=2.5cm}

\title{Guía de Ejercicios\\Derivación Numérica}
\author{}
\date{}

\begin{document}

\maketitle
\textit{ Actúa como si todo dependiera de ti, confía como si todo dependiera de Dios.}  San Ignacio de Loyola

%\vspace{0.5cm}

\section*{Indicaciones para el desarrollo y la entrega}

\begin{itemize}

\item \textbf{Formato de entrega.}  
La entrega debe realizarse exclusivamente en formato PDF,
generado a partir de un documento escrito en \LaTeX.
No se aceptarán entregas en formato Word u otros editores.

\item \textbf{Presentación.}  
Sea ordenado y claro en la redacción.  
Incluya demostraciones completas en los ejercicios teóricos
y explique adecuadamente los resultados obtenidos en los ejercicios prácticos.
Las gráficas deben estar correctamente rotuladas y comentadas.

\item \textbf{Plazo de entrega.}  
El plazo es de \textbf{dos semanas contadas a partir del \today}.  
La fecha límite de entrega es el \textbf{12+14/03/2026}. 

\item \textbf{Uso de consultas.}  
Pueden realizar consultas durante el desarrollo de la guía.
Con todo gusto los orientaré.

\item \textbf{Clase de apoyo.}  
El martes 24 utilizaremos la clase para resolver ejercicios
y discutir dificultades comunes.  
Si es necesario, podremos organizar una clase adicional optativa
para profundizar en los temas más complejos.

\item \textbf{Uso de Inteligencia Artificial.}  
Se permite el uso de herramientas de IA como apoyo.
Sin embargo, deben ser plenamente conscientes de lo que están haciendo.
La IA puede cometer errores o generar argumentos incorrectos.
Usted es responsable del contenido que entrega.


\item {\color{red}\textbf{Entrega.}}  
Solo se deben entregar $5$ ejercicios a elección. Esta guía está pensada principalmente para estudiar y practicar, no únicamente para su entrega. Sugerencia de ejercicios para entregar: $3, 8, 9, 12, 15.$ 

\end{itemize}

\it{Las matemáticas no solamente poseen la verdad, sino la suprema belleza, una belleza fría y austera, como la de la escultura, sin atractivo para la parte más débil de nuestra naturaleza... }\\ Bertrand Russell
\begin{enumerate}

\item Explique por qué en muchos problemas científicos no es posible calcular derivadas usando la definición clásica de límite.

\item Sea una malla uniforme
\[
x_i = x_0 + ih.
\]
Explique el significado de: tamaño de paso $h$,  puntos de malla, datos discretos.


\item Muestre que la aproximación
\[
f'(x) \approx \frac{f(x+h)-f(x)}{h}
\]
tiene error de orden $O(h)$. a partir de $f'$, Obtenga una fórmula para aproximar $f''(x)$ y determine el orden del error de truncamiento

\end{enumerate}

\begin{enumerate}
\setcounter{enumi}{3}

\item Sea
\[
f(x)=\sin \pi x-e^{3x}.
\]
Aproxime $f'(1)$ usando
\[
\frac{f(x+h)-f(x)}{h}
\]
con
\[
h=0.1, \qquad h=0.01.
\]
Compare con el valor exacto.

\item Escriba un programa en MATLAB que:
\begin{itemize}
\item defina $f(x)=g(x)+h(x)$ ($g,h$ colocar como input),
\item calcule la derivada usando la fórmula hacia adelante en $x=1$, con $h$ como input.
\item calcule la derivada usando la fórmula de diferencia centrada en $x=0$, con $h$ como input.
\item compare con la derivada exacta.
\end{itemize}

\item Usando expansiones de Taylor derive la fórmula de diferencia hacia adelante de orden $O(h^2)$:
\[
f'(x) \approx
\frac{-3f(x)+4f(x+h)-f(x+2h)}{2h}.
\]

\item Derive la fórmula de diferencia hacia atrás de orden $O(h^2)$.

\item Demuestre que la fórmula centrada
\[
f'(x) \approx
\frac{f(x+h)-f(x-h)}{2h}
\]
tiene error $O(h^2)$.

\item Para
\[
f(x)=\ln(x)
\]
aproxime $f'(2)$ usando diferencias centradas de orden $4$. con
\[
h=0.1 \quad \text{y} \quad h=0.01.
\]

Compare los errores.


\item Suponga
\[
D(h)=f'(x)+Ch^2+O(h^3).
\]
Demuestre que
\[
D_R=\frac{4D(h/2)-D(h)}{3}
\]
elimina el término dominante del error.

\item Calcule
\[
f'(0), \qquad f(x)=\sin x
\]
usando diferencias centradas con
\[
h=0.1 \quad \text{y} \quad h=0.05.
\]

Aplique extrapolación de Richardson.

\item Escriba un programa que:
\begin{itemize}
\item calcule $D(h)$,
\item calcule $D(h/2)$,
\item aplique Richardson,
\item compare los errores.
\end{itemize}



\item Sea
\[
f(x)=x^2
\]
con malla
\[
0,\;0.25,\;0.5,\;0.75,\;1.
\]

Calcule la derivada usando diferencias centradas y compare con la derivada exacta.

\item Escriba un programa que:
\begin{itemize}
\item construya una malla en $[0,1]$,
\item calcule la derivada con diferencias centradas,
\item grafique la derivada exacta, la derivada numérica y el error.
\end{itemize}


\item Compare los siguientes métodos para
\[
f(x)=\sin x
\]

\begin{itemize}
\item diferencia hacia adelante
\item diferencia hacia atrás
\item diferencia centrada
\item fórmula centrada de orden $O(h^4)$
\item extrapolación de Richardson
\item Grafique el error en función de $h$.
\end{itemize}

\it{Un matemático que no es también algo de poeta nunca será un matemático completo. Karl Weierstrass
 ... excelsas, supremas, excelentísimas, incomprensibles, inestimables, imnumerables, admirables, inefables, singulares ..., que corresponden por semejanza a Dios mismo.} \\

Luca Pacioli (En su tratado sobre La Divina Proporción)

\end{enumerate}

\end{document}